%-------------------------
% Resume in Latex
% Author : Sourabh Bajaj + some brand new features from Mary Feofanova
% Website: https://github.com/sb2nov/resume
% License : MIT
%------------------------

\documentclass[12pt,a4paper]{article}
\usepackage[utf8]{inputenc}
\usepackage[T2A]{fontenc}
\usepackage[russian]{babel}

\usepackage{makecell}
\usepackage[link=off]{phonenumbers}

\usepackage{latexsym}
\usepackage[empty]{fullpage}
\usepackage{titlesec}
\usepackage{marvosym}
\usepackage[usenames,dvipsnames]{color}
\usepackage{verbatim}
\usepackage{enumitem}
\usepackage[pdftex]{hyperref}
\usepackage{fancyhdr}


\pagestyle{fancy}
\fancyhf{} % clear all header and footer fields
\fancyfoot{}
\renewcommand{\headrulewidth}{0pt}
\renewcommand{\footrulewidth}{0pt}
\usepackage[margin=0.3in]{geometry}
% Adjust margins
\addtolength{\oddsidemargin}{-0.0in}
\addtolength{\evensidemargin}{-0.0in}
\addtolength{\textwidth}{0in}
\addtolength{\topmargin}{20pt}
\addtolength{\textheight}{0.0in}

\urlstyle{same}

\usepackage{xcolor}% http://ctan.org/pkg/xcolor
\usepackage{hyperref}% http://ctan.org/pkg/hyperref
\hypersetup{
  colorlinks=true,
  linkcolor=blue!50!red,
  linkbordercolor=red,
  urlcolor=blue!70!black
}

\raggedbottom
\raggedright
\setlength{\tabcolsep}{0in}

% Sections formatting
\titleformat{\section}{
  \vspace{-10pt}\scshape\raggedright\large
}{}{0em}{}[\color{black}\titlerule \vspace{5pt}]

%-------------------------
% Custom commands
\def \ifempty#1{\def\temp{#1} \ifx\temp\empty }

\newcommand{\resumeItem}[2]{
  \item\small{
  	\ifempty{#1}#2\else\textbf{#1}{: #2 \vspace{-2pt}}\fi
  }
}

\usepackage[dvipsnames]{xcolor}
\definecolor{mygray}{gray}{0}
\usepackage{fancybox}

\usepackage{lmodern}
\usepackage{tikz}

% Style definition
\tikzset{rndblock/.style={rounded corners,rectangle,draw,outer sep=0pt}}

% Command Definition
% 1 optional to customize the aspect, 2 mandatory: text to be framed
\newcommand{\tframed}[2][]{\tikz[baseline=(h.base)]\node[rndblock,#1] (h) {\color{black}{#2}};}

\newcommand*{\mystrut}{\rule[-0.2\baselineskip]{0pt}{0.8\baselineskip}}
\newcommand{\skill}[1]{\tframed[lightgray]{\mystrut#1}}


\newcommand{\resumeSubheading}[5]{
  \vspace{-1pt}\item
    \begin{tabular*}{0.97\textwidth}{l@{\extracolsep{\fill}}r}
      \textbf{#1} & \textcolor{mygray}{#2} \\
      \textit{\small#3} & \textcolor{mygray}{\textit{\small #4}} \\
      \textit{\small#5} \\
    \end{tabular*}\vspace{5pt}
}

\newcommand{\resumeExpSubheading}[5]{
  \vspace{-1pt}\item
    \begin{tabular*}{0.97\textwidth}{l@{\extracolsep{\fill}}r}
      \textbf{#1}  & \textcolor{mygray}{#2} \\
      \textit{\small#3} & \textcolor{mygray}{\textit{\small #4}} \\
      {\scriptsize#5}
    \end{tabular*}\vspace{4pt}
}

\newcommand{\resumeProjSubheading}[4]{
  \vspace{-1pt}\item
    \begin{tabular*}{0.97\textwidth}{l@{\extracolsep{\fill}}r}
      \textbf{#1}  & \textcolor{mygray}{#2} \\
      \scriptsize {#3} & \textcolor{mygray}{\textit{\small #4}} \\
    \end{tabular*}\vspace{8pt}
}

\newcommand{\resumeSubItem}[2]{\resumeItem{#1}{#2}\vspace{-8pt}}

\renewcommand{\labelitemii}{$\circ$}

\newcommand{\resumeSubHeadingListStart}{\begin{itemize}[leftmargin=*]}
\newcommand{\resumeSubHeadingListEnd}{\end{itemize}}
\newcommand{\resumeItemListStart}{\begin{itemize}[leftmargin=0.2in]}
\newcommand{\resumeItemListEnd}{\end{itemize}\vspace{-5pt}}

\usepackage{changepage}
\newcommand{\resumeDesc}[1]{\begin{adjustwidth}{5pt}{0pt}\vspace{-2pt}{\small{#1}}\end{adjustwidth}}

%-------------------------------------------
%%%%%%  CV STARTS HERE  %%%%%%%%%%%%%%%%%%%%%%%%%%%%


\begin{document}

\selectlanguage{russian}

%----------HEADING-----------------
\begin{tabular*}{\textwidth}{l@{\extracolsep{\fill}}r}
  \textbf{\Large Алина Лисянская} & Email\vspace{7pt}: \href{mailto:arlisyanskaya@edu.hse.ru}{arlisyanskaya@edu.hse.ru}\\
  
  Github: \href{https://github.com/kurabier/}{kurabier} & Телефон : +7\hspace{0.5ex}915\hspace{0.5ex}402\hspace{0.5ex}35\hspace{0.5ex}50 \\
\end{tabular*}
\vspace{9pt}

%-----------EDUCATION-----------------
\section{Образование}
  \resumeSubHeadingListStart
    \resumeSubheading
      {Национальный исследовательский вуз «Высшая школа экономики» (НИУ ВШЭ)}{Москва, Россия}
        {ПМИ, компьютерные науки и анализ данных}{Сент. 2024 - Авг. 2028}{Бакалавр, 2 курс}
    \resumeSubheading
      {лицей НИУ ВШЭ}{Москва, Россия}
        {информатика, математика и инженерия}{Сент. 2019 - Авг. 2021}{}
  \resumeSubHeadingListEnd
  
\section{Проекты}
  \resumeSubHeadingListStart
    \resumeProjSubheading
      {\href{https://kurabier.github.io/informatics_guide/informatics_guide.pdf}{Пособие для подготовки к ЕГЭ по информатике}}{}
      {\skill{Latex} \skill{Overleaf} \skill{Git} \skill{Python} \skill{GitHub}}{ февраль 2026}
          \resumeDesc{Разработала и сверстала в \textbf{LaTeX} с нуля учебное пособие для подготовки к ЕГЭ по информатике. Материал содержит теоретическую базу, ссылки на типовые задания и решения — как аналитические, так и с примерами кода на \textbf{Python}.}

\textbf{Более подробно о проекте:}
\begin{itemize}[leftmargin=*]
    \item На базе пакета \texttt{listings} создала кастомизированный стиль для Python с подсветкой синтаксиса (\texttt{xcolor}), нумерацией строк и фоновым цветом для читаемости кода;
    \item Настроила геометрию страницы (\texttt{geometry}), внедрила гиперссылки (\texttt{hyperref}) на задания. Локализовала на русский язык (\texttt{babel}) и настроила «Содержание» для документа;
    \item Вела разработку в \textbf{Linux} с использованием \textbf{Bash-скриптов} для автоматизации сборки проекта. Настроила \textbf{Git} для контроля версий и организовала удаленное хранение на \textbf{GitHub}. Реализовала \textbf{GitHub Actions workflow} для автоматической компиляции \texttt{.tex} в \texttt{.pdf} при каждом обновлении;
    \item Опубликовала готовое пособие на \textbf{GitHub Pages} — пользователи всегда имеют доступ к актуальной версии по постоянной ссылке, а любые изменения в репозитории обновляют документ автоматически.
\end{itemize}
          \resumeProjSubheading
      {\href{https://kurabier.github.io/math_guides/math_guide.pdf}{Пособие для подготовки к ЕГЭ по профильной математике}}{}
      {\skill{Latex} \skill{Overleaf} \skill{Git} \skill{Python} \skill{GitHub}}{Март 2026}
          \resumeDesc{Разработала и сверстала в \textbf{LaTeX} учебное пособие для подготовки к ЕГЭ по профильной математике, ориентированное на планиметрию и векторы.}

\textbf{Более подробно о проекте:}

\begin{itemize}[leftmargin=*]
    \item Создала геометрические иллюстрации с нуля, используя библиотеку \texttt{tikz} (включая библиотеки \texttt{trees} и \texttt{positioning}) для идеального качества изображений при печати и масштабировании. Настроила геометрию страницы (\texttt{geometry}) для удобного формата под печать (A4 с оптимизированными полями) и разбила материал на секции;
    \item Активно использовала пакет \texttt{amsmath} для безупречного набора сложных формул, векторных выражений и систем уравнений;
    \item Также использовала Linux и Git для разработки, GitHub — для хранения, GitHub Actions — для автоматической сборки.
\end{itemize}
      
  \resumeSubHeadingListEnd
%

\section{Опыт работы}
  \resumeSubHeadingListStart
    \resumeSubheading
      {Учебный ассистент по курсу «Инструменты промышленной разработки»}{Москва, Россия}
        {Национальный исследовательский вуз «Высшая школа экономики» (НИУ ВШЭ)}{Сент. 2025 - дек. 2025}{}
    \resumeSubheading
      {Репетитор по математике и информатике для старших классов}{Москва, Россия}
        {Частное преподавание}{Сент. 2024 - now}{}
    \resumeSubheading
      {Преподаватель по шахматам в общеобразовательной школе}{Санкт-Петербург, Россия}
        {«Онлайн-школа №1»}{Сент. 2022 - Авг. 2024}{}
  \resumeSubHeadingListEnd
%--------PROGRAMMING SKILLS------------
\section{Навыки}
 \resumeSubHeadingListStart
  \item{
     \textbf{Основное}{: Latex, Overleaf, Git, Python, Linux, Bash}
   \item{
     \textbf{Дополнительно}{: теория вероятностей, линейная алгебра, математический анализ, C++, SQL, matplotlib, numpy, pandas}
   }
   \vspace{0pt}
   }
 \resumeSubHeadingListEnd

%-------------------------------------------
\end{document}

